\section{Introdução}
\label{sec:intro}

A qualidade da educação básica é um dos principais determinantes do desenvolvimento humano e econômico de longo prazo \citep{hanushek2012}. No Brasil, o Índice de Desenvolvimento da Educação Básica (IDEB) monitora o ensino público combinando fluxo escolar e proficiência em língua portuguesa e matemática numa medida comparável entre municípios, estados e redes de ensino \citep{fernandes2007}. Apesar dos avanços das últimas duas décadas, as disparidades regionais persistem: estados do Norte e do Nordeste ficam sistematicamente abaixo da média nacional, reflexo de desigualdades estruturais de renda, infraestrutura e serviços públicos \citep{inep2025}.

O Maranhão está entre os casos mais críticos. Com a menor renda domiciliar \textit{per capita} do país no Censo Demográfico 2022 e indicadores educacionais historicamente abaixo da média nacional, o estado concentra municípios em alta vulnerabilidade socioeconômica \citep{ibge2022, inep2025}. Ainda assim, o desempenho varia entre municípios de condições semelhantes, o que coloca uma pergunta direta: o que diferencia, dentro de um mesmo patamar de vulnerabilidade, os municípios de maior IDEB daqueles de menor IDEB?

A literatura internacional examinou amplamente a relação entre condições socioeconômicas, infraestrutura escolar, qualificação docente e o desempenho dos estudantes \citep{rivkin2005, coleman1966}. No Brasil, a pesquisa em eficácia escolar mostrou que o nível socioeconômico dos alunos é o fator mais associado ao desempenho, mas que as escolas também fazem diferença \citep{brooke2008}. Estudos com dados do IDEB e do Censo Escolar associaram o desempenho à infraestrutura e à complexidade das escolas \citep{alves2013, vasconcelos2021}. Já a formação específica em licenciatura dos professores da educação básica permanece como meta nacional ainda não cumprida \citep{brasil2014, gatti2019}. Análises voltadas especificamente para o Maranhão que isolem fatores intracontextuais (aqueles que se mantêm associados ao desempenho mesmo entre municípios de contexto socioeconômico semelhante) são escassas.

Métodos de aprendizado de máquina têm sido aplicados a dados educacionais com resultados promissores, especialmente pela capacidade de lidar com bases administrativas grandes e heterogêneas \citep{ersozlu2024, kalita2025, romero2020}. Algoritmos de agrupamento identificam estruturas nos dados sem categorias predefinidas, o que os torna úteis para segmentar contextos antes de comparações \citep{faceli2021, kalita2025}. Quando a segmentação é seguida de testes não paramétricos aplicados dentro de cada grupo, torna-se possível comparar municípios de desempenho distinto sem confundir fatores escolares com diferenças de contexto.

O objetivo geral deste artigo é identificar, entre municípios maranhenses de contexto socioeconômico semelhante, os fatores escolares associados ao desempenho superior no IDEB 2025 dos anos iniciais do ensino fundamental da rede municipal. Duas questões de pesquisa (QP) orientam o estudo:

\begin{description}
    \item[QP1.] Em quantos grupos socioeconômicos os municípios maranhenses se organizam, considerando a renda \textit{per capita} e a pobreza extrema, e quão estável é essa partição?
    \item[QP2.] Dentro de cada grupo, quais fatores de infraestrutura, qualificação docente e fluxo escolar diferenciam os municípios de alto e de baixo IDEB?
\end{description}

Para respondê-las, a análise abrange os 216 municípios maranhenses com IDEB 2025 divulgado, de um total de 217, em dois estágios. No primeiro, os municípios são agrupados por contexto socioeconômico via \textit{K-Means}, usando a renda \textit{per capita} e a proporção de moradores em pobreza extrema, ambas do Censo Demográfico 2022. No segundo, dentro de cada grupo, o teste de Mann-Whitney, acompanhado de tamanho de efeito e de correção para comparações múltiplas, identifica quais características escolares e socioeconômicas diferem entre municípios de alto e baixo desempenho. As variáveis analisadas vêm do Censo Escolar 2025: infraestrutura (internet, biblioteca, laboratório de informática, saneamento), qualificação docente e indicadores de fluxo escolar.

O trabalho faz três contribuições principais. Primeiro, identifica os fatores escolares associados ao desempenho superior no IDEB dentro de cada perfil socioeconômico, tanto no grupo de municípios mais vulneráveis quanto no de renda intermediária, com destaque para a formação docente. Segundo, oferece evidência sobre o Maranhão, estado pouco estudado nessa literatura e no qual a comparação direta com médias nacionais é pouco informativa. Terceiro, aplica uma abordagem em dois estágios (segmentação socioeconômica seguida de análise intragrupo), construída com dados públicos e \textit{scripts} abertos, que pode ser replicada em outras unidades federativas. Os resultados oferecem subsídios a gestores e formuladores de políticas que precisam decidir onde concentrar recursos.

O restante do artigo está organizado da seguinte forma. A Seção~\ref{sec:relacionados} apresenta o referencial teórico sobre os fatores associados ao desempenho no IDEB e sobre a mineração de dados educacionais, e posiciona o estudo em relação aos trabalhos relacionados. A Seção~\ref{sec:metodos} descreve as fontes de dados, as variáveis e os dois estágios da análise. A Seção~\ref{sec:resultados} apresenta a segmentação dos municípios e os fatores discriminantes, organizados pelas questões de pesquisa e discutidos à luz da literatura. A Seção~\ref{sec:limitations} discute as limitações, e a Seção~\ref{sec:conclusion} sintetiza os achados e aponta trabalhos futuros.
